% ═══════════════════════════════════════════════════════════════════════════════
% Frozen Brane Boundary Conditions: Complete Analysis
% ═══════════════════════════════════════════════════════════════════════════════
% TITLE: Frozen on Brane via ξ-BC: Model, Derivation, and Verdict
% STATUS: [Der] - Derived from action variation and Green's function analysis
% ORIGIN: aside_frozen_brane_bc_v1/ (consolidated from 6 markdown files)
% ═══════════════════════════════════════════════════════════════════════════════

\documentclass[11pt,a4paper]{article}
\usepackage{amsmath,amssymb,amsthm}
\usepackage{booktabs}
\usepackage{enumitem}
\usepackage[margin=2.5cm]{geometry}

\newtheorem{theorem}{Theorem}
\newtheorem{lemma}{Lemma}
\newtheorem{definition}{Definition}

\title{Frozen on Brane via $\xi$-BC:\\Complete Analysis of Boundary Conditions and Vortex Interactions}
\author{EDC Derivation Library}
\date{January 2026}

\begin{document}
\maketitle

\begin{abstract}
We investigate whether ``frozen on brane'' boundary conditions in the fifth dimension $\xi$
can create an attractive regime in the vortex--vortex interaction potential. Starting from
action variation, we derive Neumann, Robin, and Dirichlet BC as limiting cases. We solve
the eigenvalue problem for each BC type and compute the effective 2D Green's function
under various $\xi$-weightings. The conclusion is definitive: $V'_{\mathrm{lin}}(d) > 0$
for all BC and weight choices. The interaction minimum arises from topological core
repulsion (radial-frozen profiles), not from $\xi$-BC effects.
\end{abstract}

\tableofcontents
\newpage

% ═══════════════════════════════════════════════════════════════════════════════
\section{Model and Definitions}
\label{sec:model}
% ═══════════════════════════════════════════════════════════════════════════════

\subsection{Geometry}

The thick brane occupies coordinates:
\begin{itemize}[nosep]
    \item \textbf{Transverse 2D:} $(r, \theta)$ --- the plane where vortices live
    \item \textbf{Fifth dimension:} $\xi \in [0, \delta]$ --- brane thickness
\end{itemize}

The metric is:
\begin{equation}
ds^2 = dr^2 + r^2 d\theta^2 + d\xi^2
\end{equation}

\subsection{Field Content}

Complex scalar field $\Phi(r, \theta, \xi)$ with vortex ansatz:
\begin{equation}
\Phi(r, \theta, \xi) = f(r, \xi) \cdot e^{in\theta}
\end{equation}
where $n \in \mathbb{Z}$ is the winding number and $f(r, \xi) \geq 0$ is the profile.

\subsection{Two Distinct ``Frozen'' Concepts}

\begin{definition}[Radial Step Frozen]
The radial profile is a step function:
\begin{equation}
f(r) = \Theta(r - a) = \begin{cases} 0, & r < a \\ v, & r \geq a \end{cases}
\end{equation}
This eliminates the GL coherence length, giving exact geometric coefficients
(e.g., $4\pi/3$, $2\pi^2$) required for $m_p/m_e = 6\pi^5$.
\end{definition}

\begin{definition}[Brane-Frozen ($\xi$-BC)]
The $\xi$-profile is pinned at boundaries:
\begin{itemize}[nosep]
    \item Neumann: $\partial_\xi f|_{\text{boundaries}} = 0$
    \item Dirichlet: $f|_{\text{boundaries}} = v$ (fixed value)
    \item Robin: $\partial_\xi f = \kappa f$ at boundaries
\end{itemize}
\end{definition}

\textbf{These are independent concepts.} Both can be true simultaneously.

% ═══════════════════════════════════════════════════════════════════════════════
\section{Boundary Conditions from Action Variation}
\label{sec:bc-derivation}
% ═══════════════════════════════════════════════════════════════════════════════

\subsection{Total Action}

\begin{equation}
S = S_{\text{bulk}} + S_{\partial,0} + S_{\partial,\delta}
\end{equation}

where the bulk action is:
\begin{equation}
S_{\text{bulk}} = 2\pi \int_0^\infty r\,dr \int_0^\delta d\xi \left[
\frac{1}{2}(\partial_r f)^2 + \frac{n^2 f^2}{2r^2} + \frac{1}{2}(\partial_\xi f)^2 + U(f)
\right]
\end{equation}

and boundary terms:
\begin{equation}
S_{\partial} = 2\pi \int r\,dr \left[ V_0(f)|_{\xi=0} + V_\delta(f)|_{\xi=\delta} \right]
\end{equation}

\subsection{Variation and Boundary Conditions}

Varying $f \to f + \delta f$ and integrating by parts in $\xi$:
\begin{equation}
\int_0^\delta d\xi\, (\partial_\xi f)(\partial_\xi \delta f) =
[(\partial_\xi f)\delta f]_0^\delta - \int_0^\delta d\xi\, (\partial_\xi^2 f)\delta f
\end{equation}

Collecting boundary terms:
\begin{align}
\text{At } \xi = 0: \quad & \partial_\xi f|_{\xi=0} = V'_0(f)|_{\xi=0} \\
\text{At } \xi = \delta: \quad & \partial_\xi f|_{\xi=\delta} = -V'_\delta(f)|_{\xi=\delta}
\end{align}

\begin{theorem}[BC from Action --- \textnormal{[Der]}]
\label{thm:bc-from-action}
The boundary conditions follow from the variational principle:
\begin{center}
\begin{tabular}{@{}lll@{}}
\toprule
Boundary Action & BC at $\xi = 0$ & BC at $\xi = \delta$ \\
\midrule
$V = 0$ (none) & $\partial_\xi f = 0$ (Neumann) & $\partial_\xi f = 0$ (Neumann) \\
$V = \frac{1}{2}\kappa f^2$ & $\partial_\xi f = \kappa f$ (Robin) & $\partial_\xi f = -\kappa f$ (Robin) \\
$V = \frac{1}{2}\lambda(f-v_0)^2$, $\lambda\to\infty$ & $f = v_0$ (Dirichlet) & $f = v_\delta$ (Dirichlet) \\
\bottomrule
\end{tabular}
\end{center}
\end{theorem}

\textbf{``Frozen on brane'' = Dirichlet BC in the stiff limit.}

% ═══════════════════════════════════════════════════════════════════════════════
\section{Mode Spectrum Under Various BC}
\label{sec:mode-spectrum}
% ═══════════════════════════════════════════════════════════════════════════════

We solve the eigenvalue problem:
\begin{equation}
-\frac{d^2\phi_m}{d\xi^2} = \mu_m^2 \phi_m \quad \text{on } \xi \in [0, \delta]
\end{equation}

\subsection{Neumann--Neumann (Baseline)}

\textbf{Eigenvalues:} $\mu_m = m\pi/\delta$, \quad $m = 0, 1, 2, \ldots$

\textbf{Eigenfunctions:} $\phi_m(\xi) = \cos(m\pi\xi/\delta)$

\textbf{Zero mode ($m=0$):} $\mu_0 = 0$, $\phi_0 = 1$ (constant). \textbf{EXISTS.}

\subsection{Dirichlet--Dirichlet}

\textbf{Eigenvalues:} $\mu_m = m\pi/\delta$, \quad $m = 1, 2, 3, \ldots$

\textbf{Eigenfunctions:} $\phi_m(\xi) = \sin(m\pi\xi/\delta)$

\textbf{Zero mode:} \textbf{DOES NOT EXIST} (excluded by BC).

\subsection{Robin--Robin (General)}

Transcendental eigenvalue equation:
\begin{equation}
\tan(\mu\delta) = \frac{\mu(\kappa_0 + \kappa_\delta)}{\mu^2 - \kappa_0\kappa_\delta}
\end{equation}

\textbf{Generic Robin: NO zero mode} ($\mu = 0$ not allowed).

\subsection{Summary Table}

\begin{center}
\begin{tabular}{@{}lccl@{}}
\toprule
BC Type & Zero Mode? & Lowest Mode & Effect on log term \\
\midrule
Neumann--Neumann & YES ($m=0$) & $\mu_0 = 0$ & Log term PRESENT \\
Dirichlet--Dirichlet & NO & $\mu_1 = \pi/\delta$ & Log term ABSENT \\
Robin--Robin (generic) & NO & $\mu_1 > 0$ & Log term ABSENT \\
Dirichlet (matching $v$) & YES$^*$ & Constant $v$ & Log term PRESENT$^*$ \\
\bottomrule
\end{tabular}
\end{center}
$^*$For Dirichlet with matching boundary values, the constant $v$ acts like a zero mode.

% ═══════════════════════════════════════════════════════════════════════════════
\section{Green's Function with Weighting}
\label{sec:greens-function}
% ═══════════════════════════════════════════════════════════════════════════════

\subsection{Effective 2D Green's Function}

\begin{equation}
G_{\text{eff}}(r) = \int_0^\delta d\xi \int_0^\delta d\xi'\, w(\xi) w(\xi') G(r; \xi, \xi')
\end{equation}

where $w(\xi)$ is a weight function with $\int_0^\delta w(\xi)\,d\xi = 1$.

Mode expansion:
\begin{equation}
G_{\text{eff}}(r) = \sum_m \frac{W_m^2}{N_m} G_m(r), \quad
W_m := \int_0^\delta w(\xi) \phi_m(\xi)\,d\xi
\end{equation}

\subsection{2D Green's Functions}

\textbf{Zero mode ($\mu_0 = 0$):}
\begin{equation}
G_0(r) = -\frac{1}{2\pi} \ln(r/L)
\end{equation}

\textbf{Massive modes ($\mu_m > 0$):}
\begin{equation}
G_m(r) = \frac{1}{2\pi} K_0(\mu_m r)
\end{equation}

\subsection{Summary: BC and Weight Combinations}

\begin{center}
\begin{tabular}{@{}llccc@{}}
\toprule
BC & Weight & Zero Mode? & Log Term? & $V'(d)$ Sign \\
\midrule
Neumann & Uniform & Yes & Yes & $\mathbf{> 0}$ \\
Neumann & Boundary & Yes & Yes & $\mathbf{> 0}$ \\
Dirichlet (hom.) & Uniform & No & No & $\mathbf{> 0}$ \\
Robin & Uniform & No & No & $\mathbf{> 0}$ \\
Dirichlet (matching $v$) & Any & Yes$^*$ & Yes$^*$ & $\mathbf{> 0}$ \\
\bottomrule
\end{tabular}
\end{center}

% ═══════════════════════════════════════════════════════════════════════════════
\section{Sign Analysis: The No-Go Result}
\label{sec:sign-analysis}
% ═══════════════════════════════════════════════════════════════════════════════

\begin{theorem}[BC and Weight Do Not Create Attraction --- \textnormal{[Der]}]
\label{thm:no-attraction}
For any combination of:
\begin{itemize}[nosep]
    \item Boundary conditions: Neumann, Robin, or Dirichlet
    \item $\xi$-weight: uniform or boundary-localized
\end{itemize}
the linearized interaction $V_{\mathrm{lin}}(d)$ for same-sign vortices satisfies:
\begin{equation}
V'_{\mathrm{lin}}(d) > 0 \quad \text{for all } d > 0
\end{equation}
\end{theorem}

\begin{proof}
Two cases:
\begin{enumerate}
    \item \textbf{Zero mode exists:} $V$ includes log term. Derivative:
    $\frac{d}{dd}[\ln d] = \frac{1}{d} > 0$.

    \item \textbf{No zero mode:} $V$ includes only $K_0$ terms with positive coefficients.
    Derivative: $\frac{d}{dd}[K_0(\mu d)] = -\mu K_1(\mu d)$, but entering $V$ with sign
    that gives $V' \propto +\mu K_1(\mu d) > 0$.
\end{enumerate}
In both cases, $V'$ is a sum of positive terms.
\end{proof}

\textbf{The no-go result is ROBUST to BC and weight changes.}

% ═══════════════════════════════════════════════════════════════════════════════
\section{Relation to Radial Step Frozen}
\label{sec:radial-comparison}
% ═══════════════════════════════════════════════════════════════════════════════

\subsection{Which Addresses the No-Go?}

\begin{center}
\begin{tabular}{@{}lcc@{}}
\toprule
Aspect & Radial-Frozen $f(r)$ & Brane-Frozen ($\xi$-BC) \\
\midrule
Definition & $f(r) = \Theta(r-a)$ & BC at $\xi = 0, \delta$ \\
Direction & Radial ($r$) & Fifth dimension ($\xi$) \\
Creates core energy? & YES & NO \\
$V_{\text{core}} \to +\infty$ at $d\to 0$? & YES (Theorem 2) & NO \\
Changes $V'_{\text{lin}}$ sign? & N/A & NO \\
Resolves no-go? & \textbf{YES} & \textbf{NO} \\
\bottomrule
\end{tabular}
\end{center}

\subsection{The Minimum Mechanism}

The minimum in $V(d)$ arises from:
\begin{itemize}[nosep]
    \item $V_{\text{core}} \to +\infty$ as $d \to 0$ \quad [from radial-frozen core, topology]
    \item $V_{\text{lin}} \to +\infty$ as $d \to \infty$ \quad [from log/Yukawa growth]
    \item Continuity $\Rightarrow$ minimum at some $d_0$
\end{itemize}

\textbf{BC provide scale $\delta$ and mode spectrum, NOT the minimum mechanism.}

% ═══════════════════════════════════════════════════════════════════════════════
\section{Verdict}
\label{sec:verdict}
% ═══════════════════════════════════════════════════════════════════════════════

\subsection{Answers to Key Questions}

\textbf{Q1: Does ``frozen on brane via BC'' exist as a derived stiff limit?}\\
\textbf{YES [Der].} Dirichlet BC emerges as $\lambda\to\infty$ limit of pinning potential.

\textbf{Q2: Does it create a minimum in $V_{\mathrm{lin}}(d)$?}\\
\textbf{NO [Der].} All BC scenarios give $V'_{\mathrm{lin}}(d) > 0$. BC affect mode spectrum but NOT sign structure.

\textbf{Q3: Does the minimum still rely on core overlap/topology?}\\
\textbf{YES [Der].} Minimum comes from radial-frozen core (Theorem 2), not BC effects.

\subsection{Epistemic Status}

\begin{center}
\begin{tabular}{@{}ll@{}}
\toprule
Result & Status \\
\midrule
BC from action variation & [Der] \\
Mode spectrum under BC & [Der] \\
$V'_{\mathrm{lin}} > 0$ for all BC & [Der] \\
$V_{\text{core}} \to +\infty$ as $d \to 0$ & [Der] (topology) \\
Minimum exists & [Dc] (continuity) \\
Minimum location $d_0 \sim \delta$ & [P] or [Cal] \\
``BC create attraction'' & \textbf{FALSE} \\
\bottomrule
\end{tabular}
\end{center}

\subsection{The Honest Statement}

\begin{quote}
``Thick-brane boundary conditions provide the characteristic scale $\delta$ and mode
structure, but do not create attraction or a minimum in $V_{\mathrm{lin}}(d)$.
The minimum in the full potential $V(d) = V_{\text{core}} + V_{\text{lin}}$ arises
from the balance of topological core repulsion [Der] and logarithmic confinement [Dc],
not from $\xi$-BC effects. Chapter 2's radial-frozen approach is validated and essential.''
\end{quote}

\end{document}
